\Article{3}
Wenn $m$ eine komplexe ganze Zahl bezeichnet, die mit $\mu$ keinen gemeinsamen Teiler hat, und sämtliche Zahlen eines Viertelsystems $\kappa$ mit $m$ multipliziert werden, so bilden die Produkte oder beliebige ihnen kongruente Zahlen ihrerseits ein Viertelsystem. Ebenso entstehen durch Multiplikation der Zahlen der Systeme $\kappa',\kappa'',\kappa'''$ drei weitere Viertelsysteme, die mit dem ersten zusammengehören. Der Beweis lässt sich leicht führen. Diese vier neuen Systeme mögen mit $m\kappa,m\kappa',m\kappa'',m\kappa'''$ bezeichnet werden, unabhängig davon, ob die Produkte selbst oder ihnen kongruente Zahlen gewählt werden. Im letzteren Fall kann man stets Zahlen wählen, die in einem der Systeme $\kappa,\kappa',\kappa'',\kappa'''$ liegen. So ist das System $\chi$, wenigstens im Allgemeinen gesprochen, auf zwei verschiedene Arten in Viertelsysteme zerlegt. Nehmen wir an, dass $m\kappa$ gemeinsam hat

\[
\begin{array}{rl}
\text{mit }\kappa\dotsfive&\lambda\text{ Zahlen}\\
\kappa'\dotsfive&\lambda'\text{ Zahlen}\\
\kappa''\dotsfive&\lambda''\text{ Zahlen}\\
\kappa'''\dotsfive&\lambda'''\text{ Zahlen}
\end{array}
\]
Dann haben auch $\kappa'$ mit $m\kappa'$, $\kappa''$ mit $m\kappa''$ und $\kappa'''$ mit $m\kappa'''$ jeweils $\lambda$ Glieder gemeinsam;
\SourceLeaf{380}
$\kappa''$ mit $m\kappa'$, $\kappa'''$ mit $m\kappa''$, $\kappa$ mit $m\kappa'''$ jeweils $\lambda'$ Glieder usw. Es sei $\varepsilon$ der kleinste Rest von $\lambda'+2\lambda''+3\lambda'''$ modulo $4$, oder $\varepsilon$ eine der vier Zahlen $0,1,2,3$, je nachdem $\lambda'+2\lambda''+3\lambda'''$ die Form $4n,4n+1,4n+2,4n+3$ hat. Ich behaupte nun, dass $\varepsilon$ von der Anordnung des Viertelsystems $\kappa$ unabhängig ist.

Um die Beweisführung zu erleichtern, verwende ich folgende Bezeichnung. $\Pi\psi$ soll die Zahl $0,1,2,3$ bezeichnen, je nachdem die durch $\mu$ nicht teilbare Zahl $\psi$ (selbst oder durch Kongruenzrepräsentation) in der Gruppe $\kappa,\kappa',\kappa'',\kappa'''$ liegt. Daraus ergeben sich unmittelbar folgende Aussagen:

I.\quad $\Pi(i\psi)\equiv1+\Pi\psi\ (\text{mod. }4)$.

II.\quad Die Zahl $i^{-\Pi\psi}\pdot\psi$ liegt selbst oder durch Kongruenzrepräsentation in der Gruppe $\kappa$.

III.\quad $\Sigma\Pi m\varphi\equiv\varepsilon,\ (\text{mod. }4)$, wenn die Summation über alle in $\kappa$ liegenden Glieder $\varphi$ erstreckt wird.

Es sei nun $k$ ein anderes Viertelsystem aus $f,f',f''$ usw., während $\kappa$ aus $\varphi,\varphi',\varphi''$ usw. besteht. Ich setze voraus, was erlaubt ist, da die \emph{Reihenfolge} der Glieder in $\kappa$ beliebig ist, dass $f$ mit $\varphi$ identisch ist oder zusammengehört, $f'$ mit $\varphi'$, $f''$ mit $\varphi''$ usw. Die mit $k$ zusammengehörigen Viertelsgruppen seien $k'(\equiv ik)$, $k''(\equiv-k)$, $k'''(\equiv-ik)$. Ferner habe die Charakteristik $\mathrm{P}$ für die Gruppen $k,k',k'',k'''$ dieselbe Bedeutung wie $\Pi$ für $\kappa,\kappa',\kappa'',\kappa'''$, sodass $\mathrm{P}\psi=0,1,2,3$, je nachdem $\psi$ zu $k,k',k'',k'''$ gehört.

Geht man also von der Zerlegung von $\chi$ in die Viertelsysteme $k,k',k'',k'''$ anstelle von $\kappa,\kappa',\kappa'',\kappa'''$ aus, so tritt an die Stelle von $\varepsilon$ der kleinste Rest von $\mathrm{P}mf+\mathrm{P}mf'+\mathrm{P}mf''+\mathrm{P}mf'''$ usw. usw. oder von $\Sigma\mathrm{P}mf$ modulo $4$, und es ist zu beweisen, dass $\Sigma\mathrm{P}mf-\Sigma\Pi m\varphi$ durch $4$ teilbar ist.

Wir schreiben diese Größe so:
\[
\begin{array}{l}
\phantom{+}\mathrm{P}mf+\mathrm{P}mf'+\mathrm{P}mf''+\mathrm{P}mf'''+\text{usw.}\\
-\mathrm{P}m\varphi-\mathrm{P}m\varphi'-\mathrm{P}m\varphi''-\mathrm{P}m\varphi'''-\\
+\mathrm{P}m\varphi+\mathrm{P}m\varphi'+\mathrm{P}m\varphi''+\mathrm{P}m\varphi'''+\\
-\Pi m\varphi-\Pi m\varphi'-\Pi m\varphi''-\Pi m\varphi'''-
\end{array}
\]

Da nach Voraussetzung $f$ und $\varphi$ kongruent sind oder zusammengehören, gilt dasselbe für $mf$ und $m\psi$, und es ist

\SourceLeaf{381}
\[
\left.\begin{aligned}
f&\equiv i^{-\mathrm{P}\varphi}\varphi\\
i^{-\mathrm{P}mf}mf&\equiv i^{-\mathrm{P}m\varphi}m\varphi
\end{aligned}\right\}\ \text{mod. }\mu
\]
Daraus folgt leicht $\mathrm{P}mf-\mathrm{P}m\varphi\equiv-\mathrm{P}\varphi\ (\text{mod. }4)$, und die Summe der beiden obersten Reihen $\equiv-\Sigma\mathrm{P}\varphi$. Da ferner $\mathrm{P}m\varphi-\Pi m\varphi\equiv\mathrm{P}(i^{-\Pi m\varphi}m\varphi)$, $m\varphi\pdot i^{-\Pi m\varphi}$ zu $\kappa$ gehört und die Gesamtheit \emph{aller} $m\varphi\pdot i^{-\Pi m\varphi}$ ohne Rücksicht auf die Reihenfolge mit allen $\varphi$ übereinstimmt, ist die Summe \emph{aller} $\mathrm{P}(m\varphi\pdot i^{-\Pi m\varphi})$ gleich der Summe aller $\mathrm{P}\varphi$. Folglich ist die Summe der dritten und vierten Reihe $\equiv\Sigma\mathrm{P}\varphi\ (\text{mod. }4)$, also die Summe aller vier Reihen $\equiv0\ (\text{mod. }4)$. Was zu beweisen war.

Da $\varepsilon$ also unabhängig von der Wahl der Viertelsysteme nur von $m$ und $\mu$ abhängt, nennen wir $\varepsilon$ den Charakter der Zahl $m$ in Bezug auf den Modul $\mu$ und bezeichnen ihn mit $\text{Ch. }m\ (\text{mod. }\mu)$. Offenbar ist dies nur eine Verallgemeinerung der Definition, die (Art. $\ldots$) für den Fall gegeben wurde, dass $\mu$ eine Primzahl ist, oder schließt sie ein.
